\section{Introduction}

Time series forecasting models map a history window of length $L$ to a future horizon. The choice of $L$ is one of the oldest decisions in the field, and one of the least examined. Classical practice selects it by autocorrelation diagnostics; the deep learning era replaced diagnostics with convention: the widely used long-term forecasting benchmark suite fixes $L = 96$ or $336$ for nearly every model and dataset \citep{dlinear, patchtst, itransformer, timesnet}, and extensions past a few hundred steps rarely help \citep{dlinear}. Foundation models pushed the window to thousands of steps \citep{chronos, timesfm, toto2}, and inference-time schemes stretch it further by downsampling \citep{sprint}. Across both eras, the implicit assumption is that history is a resource: more of it should not hurt, and probably helps.

We find that this assumption is wrong in a measurable, predictable, and exploitable way. Figure \ref{fig:sweep} shows the basic phenomenon. Holding the number of training windows exactly constant across lookback lengths, the test error of a standard forecaster traces a U-shape whose minimum moves by a factor of 30 across datasets: 96 steps on exchange rates, 336 on weather and minute-level ETT, 720 on ETTh1, and still improving at 2880 on electricity and traffic. The same ordering at the two poles (shortest on exchange, longest on traffic and electricity) appears in a zero-shot foundation model whose training never saw the decisive benchmark: traffic is confirmed absent from its pretraining corpus, and the one corpus overlap we find, electricity, is disclosed in Section \ref{sec:zero}. This rules out explanations based on optimization dynamics or overfitting to a training split. Lookback behaves like a resource with a dataset-specific point of diminishing, then negative, returns. We call this point the \emph{effective history span} (EHS) of a forecasting problem.

Why should distant history ever hurt? The information-theoretic intuition that more context cannot reduce a sufficient model's performance fails on time series for a structural reason: the data-generating process drifts. Windows drawn from an old regime carry an outdated conditional mapping, and empirical risk minimization mixes the old mapping with the new one. The pollution account, however, is only half of the story. We show that a single-axis theory, in which the useful history is the time since the last regime change, makes a sharp prediction---that regime-switching data should demand short windows---which is cleanly falsified: electricity and traffic switch regimes most frequently among our benchmarks, yet benefit from the longest lookbacks (Section \ref{sec:bocpd}). The resolution is a two-axis account. Non-stationarity sets the cost of distant history; periodic and long-memory structure sets its benefit, because such structure is invariant across regimes. Electricity's weekly cycle survives every regime change, while the level of a drifting exchange rate does not. The effective history span is the balance of the two axes.

This account has four consequences that we develop into results. First, the harm axis obeys a quantitative law: on controlled synthetic drift the best span scales as $W^* \propto \lambda^{-1/3}$ in the drift rate $\lambda$, refuting the classical $-2/3$ exponent (Section \ref{sec:powerlaw}). Second, the inability of standard architectures to exploit long context is mechanical and inspectable: attention weights flatten with context, no sink forms to absorb stale mass, linear models do not learn to zero distant lags, and the projection head keeps reading the full window under masking (Section \ref{sec:mechanism}). Third, the EHS is partially predictable before any training: simple context statistics, a segment-mean drift index and long-lag autocorrelation, achieve Spearman $|\rho| = 0.83$--$0.94$ with the measured long-context benefit in a six-dataset pilot (Sections \ref{sec:predict} and \ref{sec:predictfull}). Fourth, if the diagnosis is right, then allocating history like a budget should recover the loss: an input budget recovers 74--98\% of the degradation on trained patch-family models, and a statistics-guided budget on a frozen foundation model matches validation-based span selection on all eight benchmarks tested (Sections \ref{sec:predictfull} and \ref{sec:method}).

\paragraph{Contributions.}
\begin{itemize}
    \item We define the effective history span and provide a confound-controlled measurement protocol that separates window length from training data quantity (Section \ref{sec:measure}).
    \item We establish the two-axis account of when distant history helps or hurts, falsifying the single-axis changepoint alternative (Section \ref{sec:why}).
    \item We show that the EHS is partially predictable from cheap, leakage-free context statistics---the only zero-cost selector that identifies long-span datasets, where PACF/AIC order selection fail structurally---and, for a frozen foundation model, a statistic rule matches validation-based span selection at zero cost (Section \ref{sec:predictfull}).
    \item We diagnose why architectures fail to exploit long context, measure a power law $W^* \propto \lambda^{-1/3}$ for the harm axis, and introduce history budgeting, an input-side module allocating history per window (Sections \ref{sec:mechanism}--\ref{sec:method}).
\end{itemize}
